\documentclass{article}
% Source: C:/Users/rober/Downloads/Topic_02_Teacher_Edition.pdf, PDF pages 97-98 (printed pages 115A-115B): Topic Assessment Form A with the TE's marked answers, Form B with marked answers, Item Analysis for both forms. Cross-checked against the Realize download aga_24_a2_na_02_tta.docx (Form A, as text; math dropped by the conversion).
\usepackage{amsmath}
\usepackage{amssymb}
\newenvironment{assessment-meta}{}{}   % topic, title, forms, item count, printed pages, sources
\newenvironment{item-analysis}{}{}     % Savvas's Item / DOK / Skills Review and Practice table (same for both forms)
\newenvironment{ta-item}[2]{}{}        % #1 = item number, #2 = form (A or B)
\newcommand{\answer}[1]{}              % the TE's marked answer, inside a ta-item
\newcommand{\placeholder}[2]{}         % #1 = type (graph, table, figure, number-line), #2 = complete description
\newcommand{\uncertain}[2]{#1}         % #1 = best reading, #2 = alternatives and why
\newcommand{\te}[1]{}                  % transcriber's note (a discrepancy, a publisher misprint, a computed check)
\begin{document}
\begin{assessment-meta}
topic: 2
title: Quadratic Functions and Equations
forms: A and B (parallel item for item)
items: 17 per form
printed: 115A-115B
sources: C:/Users/rober/Downloads/Topic_02_Teacher_Edition.pdf, PDF pages 97-98 (Form A student pages 1 of 2 and 2 of 2 on page 97; Form B student pages 1 of 2 and 2 of 2 plus the Item Analysis on page 98); Realize download aga_24_a2_na_02_tta.docx (Form A text, cross-check only)
\end{assessment-meta}
\te{Page 98 running head reads "TOPIC 1 Solving Equations and Inequalities"; this is a publisher misprint. The Form B student page header reads "2 Topic Assessment Form B" and the TE footer reads "TOPIC 2 | 115B | TOPIC ASSESSMENT", which are authoritative. Page 97 running head reads "TOPIC 2 Quadratic Functions and Equations", footer "TOPIC 2 | 115A | TOPIC ASSESSMENT".}
\te{Page 97 TE note under the student pages: "Topic Assessment Assess students' understanding of topic concepts and skills using the Topic Assessment found at SavvasRealize.com. These auto-scored online assessments provide students with a breadth of technology-enhanced item types. There are two versions of the Topic Assessment, Form A and Form B. These two versions, available in print and at SavvasRealize.com, are parallel tests that assess the same content item for item. The Item Analysis chart on the next page can be used for both versions."}
% Source page 98: 115B Item Analysis
\begin{item-analysis}
Item & DOK & Skills Review and Practice \\
1 & 2 & F17 \\
2 & 1 & F68 \\
3 & 2 & F60 \\
4 & 2 & A29 \\
5 & 2 & F16 \\
6 & 2 & F16 \\
7 & 2 & F16 \\
8 & 2 & NS11 \\
9 & 1 & NS11 \\
10 & 2 & NS11 \\
11 & 2 & F68 \\
12 & 2 & NS11 \\
13 & 2 & F60 \\
14 & 2 & F50 \\
15 & 2 & A32 \\
16 & 2 & A34, A56 \\
17 & 2 & A34, A56 \\
\end{item-analysis}
% Source page 97: 115A Topic Assessment Form A
\begin{ta-item}{1}{A}
The distance of a golf ball from the hole can be represented by the right side of a parabola with vertex ((-1,8)). The ball reaches the hole 1 second after it is hit, at time 0. What is the equation of the parabola, in vertex form, that represents the ball's distance from the hole, (y), at time (x)?
(y=[ ](x+[ ])^2+[ ])
\answer{1A. (y=-2(x+1)^2+8): boxes -2, 1, 8}
\te{computed: (-2(1+1)^2+8=0), so the ball is at the hole at (x=1); the key is consistent with the printed prompt.}
\te{The Realize download words item 1 differently: "The graph of a golf ball's distance y from the hole over time x is a parabola, with x in seconds and y in feet. After 1 second the ball is at its maximum distance from the hole at 8 feet. After 3 seconds the ball reaches the hole. What is the equation in vertex form of the parabola?" (that version gives vertex (1,8) and (y=-2(x-1)^2+8)). The TE image rules; the printed TE wording is transcribed above.}
\end{ta-item}
\begin{ta-item}{2}{A}
Function (g) is a transformation of the parent function (f(x)=x^2). The graph of (f) is reflected across the (x)-axis, and then translated left 4 units and down 2 units to form the graph of (g). Write the equation for (g) in the form (y=ax^2+bx+c).
A. (y=-x^2+8x+14)
B. (y=-x^2-8x-18)
C. (y=x^2-8x+18)
D. (y=-x^2-8x+14)
\answer{2A. B}
\end{ta-item}
\begin{ta-item}{3}{A}
The shape of the inside of a glass follows a parabola with the function (f(x)=x^2+6x+9). What point represents the bottom of the inside of the glass?
A. ((-3,0))
B. ((0,-3))
C. ((0,3))
D. ((3,0))
\answer{3A. A}
\end{ta-item}
\begin{ta-item}{4}{A}
A rectangular painting has length 1 ft longer than its width. What is the width of the painting if its area is 12 (ft^2)?
A. 4 feet
B. 2 feet
C. 6 feet
D. 3 feet
\answer{4A. D}
\end{ta-item}
\begin{ta-item}{5}{A}
The equation (y=-4.9x^2+14) represents the height (y) in meters of a rock dropped off a bridge over time (x) in seconds. Graph this equation from the time the rock dropped to when it reached the water.
\placeholder{graph}{Answer graph (the magenta curve and dots are the TE's printed answer; the grid is the prompt). First-quadrant grid; vertical axis labeled y with an up arrow, labels 0, 4, 8, 12, 16, gridlines every 2, grid top at 18; horizontal axis labeled t (not x) with a right arrow, labels 0, 0.5, 1, 1.5, 2, 2.5, 3, gridlines every 0.25, grid right edge just past 3. Magenta curve: the right half of a downward parabola from a closed dot at (0,14) (the vertex, on the y-axis) curving down to a closed dot on the t-axis at about (1.7,0).}
\answer{5A. graph drawn: the curve from (0,14) down to about (1.69,0), closed dots at both ends}
\te{computed: (-4.9x^2+14=0) gives (x=\sqrt{14/4.9}\approx 1.69). The grid labels the horizontal axis t although the prompt names the variable x.}
\end{ta-item}
\begin{ta-item}{6}{A}
The height of a projectile launched from a 16-ft-tall tower is modeled by the equation (y=-16t^2+64t+16). Graph the equation. What is the maximum height, in feet, reached by the projectile?
\placeholder{graph}{Answer graph (the magenta curve and dots are the TE's printed answer; the grid is the prompt). First-quadrant grid; vertical axis labeled y with an up arrow, labels 0, 20, 40, 60, 80, gridlines every 10, grid top at 100; horizontal axis labeled t with a right arrow, labels 0, 1, 2, 3, 4, gridlines every 0.5, grid right edge at about 5. Magenta curve: downward parabola from a closed dot at (0,16) up to the vertex (2,80) and down to a closed dot on the t-axis at about (4.2,0).}
[ ] ft
\answer{6A. 80 ft; graph drawn: the curve from (0,16) through the vertex (2,80) to about (4.24,0), closed dots at both ends}
\end{ta-item}
\begin{ta-item}{7}{A}
A ball is thrown from the top row of seats in a stadium. The function (h(t)=-16t^2+64t+80) gives the height, in feet, of the ball (t) seconds after it is thrown. How long will it be before the ball hits the ground?
[ ] s
\answer{7A. 5 s}
\end{ta-item}
\begin{ta-item}{8}{A}
Use square roots to solve the equation (x^2=-25) over the complex numbers. Select all solutions for the equation.
A. (-5i)
B. (-5)
C. (-i)
D. (i)
E. 5
F. (5i)
\answer{8A. A, F}
\end{ta-item}
\begin{ta-item}{9}{A}
Write the product ((4+i)(4-i)) in the form (a+bi).
A. (16-i)
B. (16-i^2)
C. 17
D. 8
\answer{9A. C}
\end{ta-item}
\begin{ta-item}{10}{A}
Solve (5=x^2-10x+35) by completing the square.
A. (x=5+i) and (x=5-i)
B. (x=5+i\sqrt{5}) and (x=5-i\sqrt{5})
C. (x=-6-i\sqrt{6}) and (x=-5+i\sqrt{6})
D. (x=-5-i\sqrt{5}) and (x=-5+i\sqrt{5})
\answer{10A. B}
\te{Choice C is printed with -6 in the first solution and -5 in the second (a non-conjugate pair); transcribed as printed, likely a publisher misprint of -5. It does not affect the key.}
\end{ta-item}
\begin{ta-item}{11}{A}
A function is given by (f(x)=x^2+3x+1). Select all the statements that are true about this function.
A. The function written in vertex form is (f(x)=\left(x+\frac{3}{2}\right)^2-\frac{5}{4}).
B. The function written in vertex form is (f(x)=\left(x+\frac{5}{4}\right)^2-\frac{3}{2}).
C. The graph of the function has a minimum of (y=-\frac{5}{4}) at (x=-\frac{3}{2}).
D. The domain of the function is all real numbers.
E. The range of the function is all real numbers.
\answer{11A. A, C, D}
\end{ta-item}
\begin{ta-item}{12}{A}
Solve (x^2+3x=-4) using the Quadratic Formula.
(x=\frac{[ ]\pm i\sqrt{[ ]}}{[ ]})
\answer{12A. (x=\frac{-3\pm i\sqrt{7}}{2}): boxes -3, 7, 2}
\end{ta-item}
\begin{ta-item}{13}{A}
What is an equation for the quadratic function through the origin, ((4,2)), and ((-4,2))?
(y=\frac{[ ]}{[ ]}x^2)
\answer{13A. (y=\frac{1}{8}x^2): boxes 1, 8}
\end{ta-item}
\begin{ta-item}{14}{A}
A toy cannon ball is launched from a cannon on top of a platform. The function (h(t)=-5t^2+20t+4) gives the height, in meters, of the ball (t) seconds after it is launched. Write and solve an inequality to find the times where the ball is more than 12 meters above the ground. Round to the nearest hundredth.
(0<[ ]t^2+[ ]t+[ ])
[ ] seconds (<t<) [ ] seconds
\answer{14A. (0<-5t^2+20t+(-8)): boxes -5, 20, -8; 0.45 seconds (<t<) 3.55 seconds}
\end{ta-item}
\begin{ta-item}{15}{A}
What value(s) of (b) will cause (4x^2+bx+25=0) to have one real solution?
A. -200
B. -20
C. -2
D. 2
E. 20
F. 200
\answer{15A. B, E}
\end{ta-item}
\begin{ta-item}{16}{A}
Determine the number of real solutions of the system (\begin{cases} y=x^2+8 \\ y=x+15 \end{cases}).
A. 0
B. 1
C. 2
D. 3
\answer{16A. C}
\end{ta-item}
\begin{ta-item}{17}{A}
Solve (-3x^2+2x+4=-x-3) by writing a linear-quadratic system and using the intersection feature of a graphing calculator. Round to the nearest hundredth.
The graphing calculator shows the curve and line intersect at
(x\approx) [ ] and (x\approx) [ ].
\answer{17A. (x\approx -1.11) and (x\approx 2.11)}
\end{ta-item}
% Source page 98: 115B Topic Assessment Form B
\begin{ta-item}{1}{B}
The graph of a toy car's speed (y) over time (x) is a parabola that shows a minimum speed of 2 m/s after 3 seconds. After 5 seconds, the car's speed is 3 m/s. What is the equation in vertex form of the parabola?
(y=[ ](x+[ ])^2+[ ])
\answer{1B. (y=0.25(x+(-3))^2+2): boxes 0.25, -3, 2}
\end{ta-item}
\begin{ta-item}{2}{B}
Function (g) is a transformation of the parent function (f(x)=x^2). The graph of (f) is translated right 2 units and up 3 units, and then reflected across the (y)-axis to form the graph of (g). Write the equation for (g) in the form (y=ax^2+bx+c).
A. (y=x^2+2x+3)
B. (y=x^2+6x+7)
C. (y=x^2-4x+1)
D. (y=x^2+4x+7)
\answer{2B. D}
\end{ta-item}
\begin{ta-item}{3}{B}
The shape of an underground pool follows a parabola with the function (f(x)=x^2+8x). What point represents the bottom of the pool?
A. ((-4,16))
B. ((0,8))
C. ((-4,-16))
D. ((0,-8))
\answer{3B. C}
\end{ta-item}
\begin{ta-item}{4}{B}
A rectangular frame has length 7 in. longer than its width. What is the length of the painting if its area is 18 (in^2)?
A. 9 inches
B. 3 inches
C. 2 inches
D. 6 inches
\answer{4B. A}
\te{The prompt says "frame" then "the length of the painting"; transcribed as printed. computed: (w(w+7)=18) gives width 2, length 9, matching the key.}
\end{ta-item}
\begin{ta-item}{5}{B}
An equation that represents the height (y), in meters, of a bird over time (x), in seconds, as it lands is (y=\frac{1}{2}(x-4)^2). Graph this function from time 0 to when the bird lands.
\placeholder{graph}{Answer graph (the magenta curve, dots and the magenta axis labels are the TE's printed answer overlay; the grid is the prompt). Grid with the origin labeled O near the lower left; vertical axis labeled y with arrows at both ends, labels 2, 4, 6, 8, gridlines every 1, grid from -1 to 9; horizontal axis labeled x with arrows at both ends, labels 1, 2, 3, 4 and the origin O, gridlines every 0.5, grid from -0.5 to about 4.5; all of these labels and the axis letters x and y are printed in magenta like the curve (the grid itself is gray). Magenta curve: the left half of an upward parabola from a closed dot at (0,8) on the y-axis curving down to a closed dot at the vertex (4,0) on the x-axis.}
\answer{5B. graph drawn: the curve from (0,8) down to the vertex (4,0), closed dots at both ends}
\end{ta-item}
\begin{ta-item}{6}{B}
The height of a projectile launched from a 20-ft-tall tower is modeled by the equation (y=-16t^2+32t+20). Graph the equation. What is the maximum height, in feet, reached by the projectile?
\placeholder{graph}{Answer graph (the magenta curve and dots are the TE's printed answer; the grid is the prompt). First-quadrant grid; vertical axis labeled y with an up arrow, labels 0, 8, 16, 24, 32, gridlines every 4, grid top at 40; horizontal axis labeled t with a right arrow, labels 0, 1, 2, 3, 4, gridlines every 0.5, grid right edge at 5. Magenta curve: downward parabola from a closed dot at (0,20) up to the vertex (1,36) and down to a closed dot on the t-axis at (2.5,0).}
[ ] ft
\answer{6B. 36 ft; graph drawn: the curve from (0,20) through the vertex (1,36) to (2.5,0), closed dots at both ends}
\end{ta-item}
\begin{ta-item}{7}{B}
A ball is thrown from the top row of seats in a stadium. The function (h(t)=-16t^2+80t+96) gives the height, in feet, of the ball (t) seconds after it is thrown. How long will it be before the ball hits the ground?
[ ] s
\answer{7B. 6 s}
\end{ta-item}
\begin{ta-item}{8}{B}
Use square roots to solve the equation (x^2=-121) over the complex numbers. Select all solutions for the equation.
A. (i)
B. 11
C. (11i)
D. (-11i)
E. -11
F. (-i)
\answer{8B. C, D}
\end{ta-item}
\begin{ta-item}{9}{B}
Write the product (3i(5-3i)) in the form (a+bi).
A. (6+8i)
B. (-9+15i)
C. -6
D. (9+15i)
\answer{9B. D}
\end{ta-item}
\begin{ta-item}{10}{B}
Solve (-2=x^2+2x+8) by completing the square.
A. (x=-1-9i) and (x=-1+9i)
B. (x=-1+3i) and (x=-1-3i)
C. (x=4) and (x=2)
D. (x=\sqrt{10}-1) and (x=-\sqrt{10}-1)
\answer{10B. B}
\end{ta-item}
\begin{ta-item}{11}{B}
A function is given by (f(x)=x^2+5x+4). Select all the statements that are true about this function.
A. The graph of the function has a minimum of (y=-\frac{9}{4}) at (x=-\frac{5}{2}).
B. The function written in vertex form is (f(x)=\left(x+\frac{5}{2}\right)^2-\frac{9}{4}).
C. The graph of the function has a minimum of (y=-\frac{9}{4}) at (x=\frac{5}{4}).
D. The range of the function is all real numbers.
E. The domain of the function is all real numbers.
\answer{11B. A, B, E}
\end{ta-item}
\begin{ta-item}{12}{B}
Solve (x^2+5x=-8) using the Quadratic Formula.
(x=\frac{[ ]\pm i\sqrt{[ ]}}{[ ]})
\answer{12B. (x=\frac{-5\pm i\sqrt{7}}{2}): boxes -5, 7, 2}
\end{ta-item}
\begin{ta-item}{13}{B}
What is an equation for the quadratic function through the origin, ((10,2.5)), and ((-10,2.5))?
(y=\frac{[ ]}{[ ]}x^2)
\answer{13B. (y=\frac{1}{40}x^2): boxes 1, 40}
\end{ta-item}
\begin{ta-item}{14}{B}
A toy cannon ball is launched from a cannon on top of a platform. The function (h(t)=-5t^2+15t+10) gives the height, in meters, of the ball (t) seconds after it is launched. Write and solve an inequality to find the times where the ball is more than 21 meters above the ground. Round to the nearest hundredth.
(0<[ ]t^2+[ ]t+[ ])
[ ] seconds (<t<) [ ] seconds
\answer{14B. (0<-5t^2+15t+(-11)): boxes -5, 15, -11; 1.28 seconds (<t<) 1.72 seconds}
\end{ta-item}
\begin{ta-item}{15}{B}
What value(s) of (b) will cause (8x^2+bx+2=0) to have one real solution?
A. -16
B. -8
C. -4
D. 4
E. 8
F. 16
\answer{15B. B, E}
\end{ta-item}
\begin{ta-item}{16}{B}
Determine the number of real solutions of the system (\begin{cases} y=x^2+1 \\ y=1 \end{cases})
A. 0
B. 1
C. 2
D. 3
\answer{16B. B}
\end{ta-item}
\begin{ta-item}{17}{B}
Solve (-3x^2-6x+21=6x+8) by writing a linear-quadratic system and using the intersection feature of a graphing calculator. Round to the nearest hundredth.
The graphing calculator shows the curve and line intersect at
(x\approx) [ ] and (x\approx) [ ].
\answer{17B. (x\approx 0.89) and (x\approx -4.89)}
\end{ta-item}
\end{document}
